\begin{figure}[htbp]
\centering
\begin{minipage}[t]{.44\textwidth}\centering
\begin{tikzpicture}[causal,x=1cm,y=1cm,
 n/.style={draw=blue!55!black,fill=white,minimum width=8mm,minimum height=6mm,inner sep=2pt,font=\small}]
\node[n] (c10a1) at (0,2.5) {$X_1$};
\node[n] (c10a2) at (2.6,2.5) {$X_2$};
\node[n] (c10a3) at (1.3,1.2) {$X_3$};
\node[n] (c10a4) at (0,0) {$X_4$};
\node[n] (c10a5) at (2.6,0) {$X_5$};
\draw[->,green!55!black] (c10a1) -- (c10a3);
\draw[->,green!55!black] (c10a2) -- (c10a3);
\draw[->,green!55!black] (c10a3) -- (c10a4);
\draw[->,green!55!black] (c10a3) -- (c10a5);
\node[anchor=west,font=\small] at (-.5,-.6) {(a)};
\end{tikzpicture}
\end{minipage}\hfill
\begin{minipage}[t]{.52\textwidth}\centering
\begin{tikzpicture}[causal,x=1cm,y=1cm,
 n/.style={draw=blue!55!black,fill=white,minimum width=8mm,minimum height=6mm,inner sep=2pt,font=\small}]
\node[n] (c10b1) at (1.6,2.5) {$X_1$};
\node[n] (c10b2) at (0,1.2) {$X_2$};
\node[n] (c10b3) at (3.2,1.2) {$X_3$};
\node[n] (c10b4) at (1.6,0) {$X_4$};
\draw[->,green!55!black] (c10b1) -- (c10b2);
\draw[->,green!55!black] (c10b1) -- (c10b3);
\draw[->,green!55!black] (c10b2) -- (c10b4);
\draw[->,green!55!black] (c10b3) -- (c10b4);
\node[anchor=west,font=\small] at (-.5,-.6) {(b)};
\end{tikzpicture}
\end{minipage}
\par\vspace{8mm}
\begin{minipage}[t]{.40\textwidth}\centering
\begin{tikzpicture}[causal,x=1cm,y=1cm,
 n/.style={draw=blue!55!black,fill=white,minimum width=8mm,minimum height=6mm,inner sep=2pt,font=\small}]
\node[n] (c10c1) at (0,2) {$X_1$};
\node[n] (c10c2) at (1.5,2) {$X_2$};
\node[n] (c10c3) at (3,2) {$X_3$};
\node[n] (c10c4) at (0,0) {$X_4$};
\node[n] (c10c5) at (1.5,0) {$X_5$};
\node[n] (c10c6) at (3,0) {$X_6$};
\draw[<->,blue!65!black] (c10c1.north) to[bend left=35] (c10c2.north);
\draw[<->,blue!65!black] (c10c2.north) to[bend left=35] (c10c3.north);
\draw[<->,blue!65!black] (c10c1.north) to[bend left=40] (c10c3.north);
\draw[->,green!55!black] (c10c1) -- (c10c4);
\draw[->,green!55!black] (c10c2) -- (c10c5);
\draw[->,green!55!black] (c10c3) -- (c10c6);
\node[anchor=west,font=\small] at (-.5,-.6) {(c)};
\end{tikzpicture}
\end{minipage}\hfill
\begin{minipage}[t]{.58\textwidth}\centering
\begin{tikzpicture}[causal,x=1.05cm,y=1cm,
 n/.style={draw=blue!55!black,fill=white,minimum width=8mm,minimum height=6mm,inner sep=2pt,font=\small},
 lab/.style={font=\footnotesize,fill=white,inner sep=1pt}]
\node[n] (c10d1) at (2,3.2) {$X_1$};
\node[n] (c10d2) at (0,2.3) {$X_2$};
\node[n] (c10d3) at (4,2.3) {$X_3$};
\node[n] (c10d4) at (0,0) {$X_4$};
\node[n] (c10d5) at (2,-.4) {$X_5$};
\node[n] (c10d6) at (4,.4) {$X_6$};
\draw[->,red!75!black] (c10d1) -- node[lab,above left] {$b_{21}$} (c10d2);
\draw[->,red!75!black] (c10d1) -- node[lab,above right] {$b_{31}$} (c10d3);
\draw[->,red!75!black] (c10d2) -- node[lab,above,pos=.7] {$b_{32}$} (c10d3);
\draw[->,green!55!black] (c10d1) -- node[lab,left,pos=.7] {$b_{41}$} (c10d4);
\draw[->,green!55!black] (c10d1) -- node[lab,left,pos=.7] {$b_{51}$} (c10d5);
\draw[->,green!55!black] (c10d1) -- node[lab,right,pos=.4] {$b_{61}$} (c10d6);
\draw[->,green!55!black] (c10d2) -- node[lab,above,pos=.55] {$b_{26}$} (c10d6);
\draw[->,green!55!black] (c10d3) -- node[lab,right] {$b_{63}$} (c10d6);
\draw[->,red!75!black] (c10d2) -- node[lab,left,pos=.78] {$b_{25}$} (c10d5);
\draw[->,red!75!black] (c10d3) -- node[lab,right,pos=.78] {$b_{53}$} (c10d5);
\draw[->,red!75!black] (c10d4) -- node[lab,below] {$b_{54}$} (c10d5);
\draw[<->,red!75!black] (c10d2.west) to[bend right=40] node[lab,left] {$r_{24}$} (c10d4.west);
\node[anchor=west,font=\small] at (-.6,-1) {(d)};
\end{tikzpicture}
\end{minipage}
\origfig{10}{Wright 路径追踪规则示意。(a) 同一变量不能在路径中多次出现；(b) 不能先沿箭头方向前进、再逆箭头方向后退；(c) 一条路径最多只能包含一条双向边；(d) 例 1.3 的因果图。}
\label{fig:ch01c-10}
\end{figure}
